\documentclass[12pt,a4paper]{article}
\usepackage{../../../myconfig}


% Bảng trống dùng ở phần xác suất có điều kiện.
\newcommand{\bangtrong}{%
    \begin{center}
        \renewcommand{\arraystretch}{1.6}
        \begin{tabular}{|c|p{1.7cm}|p{1.7cm}|c|}
            \hline
            & $B$ & $\overline B$ & Tổng \\
            \hline
            $A$ & & & \\
            \hline
            $\overline A$ & & & \\
            \hline
            Tổng & & & $1$ \\
            \hline
        \end{tabular}
    \end{center}
}

\begin{document}

% =========================================================================
% TRANG 1: NỀN TẢNG XÁC SUẤT VÀ NGÔN NGỮ BIẾN CỐ
% =========================================================================
\titlebox{Xác suất}{Xác suất cơ bản \\ Xác suất cổ điển}{Xác suất: Xác suất cơ bản và xác suất cổ điển}

\begin{theorybox}
    \textbf{I. NỀN TẢNG}
    \[
        \underbrace{\Omega}_{\text{Mọi kết quả}}\qquad\supset\qquad\underbrace{A}_{\text{Kết quả thuận lợi}}
    \]
    \[
        \boxed{P(A)=\frac{n(A)}{n(\Omega)}}\qquad\text{(các kết quả đồng khả năng).}
    \]
\end{theorybox}

% --- Câu 1 ---
\questionLinesManual{4}{1}{%
    Gieo một xúc xắc cân đối.
    Gọi $A$: ``xuất hiện mặt chẵn''.
    Hãy viết $\Omega$, $A$ và tính $P(A)$.
}{}

\begin{theorybox}
    \textbf{II. ĐỌC ĐÚNG BIẾN CỐ}
    \[
        \begin{array}{c|c@{\qquad}c|c}
            \overline A&\text{Không }A&A\cap B&\text{Cả hai}\\[5pt]
            A\cup B&\text{Ít nhất một}&A\cap\overline B&\text{Chỉ }A
        \end{array}
    \]
\end{theorybox}

% --- Câu 2 ---
\questionLinesManual{3}{2}{%
    Lấy một bi từ mỗi hộp.
    Gọi $A$: ``hộp I lấy đỏ'', $B$: ``hộp II lấy đỏ''.
    Viết biến cố: cả hai đỏ; ít nhất một đỏ; đúng một đỏ; không có bi đỏ.
}{}

\newpage

% =========================================================================
% TRANG 2: DẠNG 1
% =========================================================================
\dangbox{1}{Ít nhất một và đúng một: khác nhau ở đâu?}

\begin{theorybox}
    \begin{minipage}{.58\linewidth}
        \[
            \boxed{P(A\cup B)=P(A)+P(B)-P(A\cap B)}
        \]
        \[
            \boxed{P(\text{đúng một})=P(A)+P(B)-2P(A\cap B)}
        \]
    \end{minipage}\hfill
    \begin{minipage}{.39\linewidth}
        \centering
        \begin{tikzpicture}[scale=.7]
            \draw[deepblue,thick] (-2,-1.3) rectangle (3,1.3);
            \fill[goldyellow!25] (0,0) circle(1);
            \fill[goldyellow!25] (1.2,0) circle(1);
            \begin{scope}
                \clip (0,0) circle(1);
                \fill[amberorange!65] (1.2,0) circle(1);

            \end{scope}
            \draw[deepblue,thick] (0,0) circle(1);
            \draw[deepblue,thick] (1.2,0) circle(1);
            \node at (-.45,0){$A$};
            \node at (1.65,0){$B$};
            \node[below] at (.6,-1.35){$A\cap B$: phần chung};
        \end{tikzpicture}
    \end{minipage}
\end{theorybox}

% --- Câu 3 ---
\questionLinesManual{7}{3}{%
    Một công ty có xác suất thắng dự án 1 là $0{,}5$, thắng dự án 2 là $0{,}6$, thắng cả hai là $0{,}4$.
    Tính xác suất thắng ít nhất một dự án và xác suất thắng đúng một dự án.
}{}

% --- Câu 4 ---
\questionLinesManual{7}{4}{%
    Cho $P(A)=0{,}6$, $P(B)=0{,}5$, $P(A\cap B)=0{,}3$.
    Tính xác suất không có biến cố nào xảy ra và xác suất chỉ có $A$ xảy ra.
}{}

\newpage

% =========================================================================
% TRANG 3: DẠNG 2
% =========================================================================
\dangbox{2}{Ít nhất một: nghĩ đến biến cố đối}

\begin{theorybox}
    \[
        \boxed{P(\overline A)=1-P(A)}
    \]
    \[
        \text{Ít nhất một}\quad\overset{\text{đối}}{\longleftrightarrow}\quad\text{Không có lần nào}.
    \]
    \[
        A,B\text{ độc lập}\ \Rightarrow\ \boxed{P(A\cap B)=P(A)P(B)}.
    \]
    \[
        \boxed{P(A\cup B)=1-P(\overline A)P(\overline B)}\quad\text{(khi độc lập).}
    \]
\end{theorybox}

% --- Câu 5 ---
\questionLinesManual{7}{5}{%
    Hai thí nghiệm độc lập, mỗi thí nghiệm thành công với xác suất $0{,}8$.
    Tính xác suất có ít nhất một thí nghiệm thành công.
}{}

% --- Câu 6 ---
\questionLinesManual{7}{6}{%
    Hai người gọi điện đều quên chữ số cuối.
    Mỗi người thử ngẫu nhiên từ 0 đến 9, không lặp số đã thử; hai người thử độc lập.
    Tính xác suất ít nhất một người gọi đúng mà không phải thử quá hai lần.
}{}

\newpage

% =========================================================================
% TRANG 4: DẠNG 3
% =========================================================================
\dangbox{3}{Biết rằng: đang xét trong tập nào?}

\begin{theorybox}
    \begin{minipage}{.6\linewidth}
        \[
            \boxed{P(A\mid B)=\frac{P(A\cap B)}{P(B)}}\quad[P(B)>0].
        \]
        \[
            \underbrace{A}_{\text{Cần tìm}}\ \big|\ \underbrace{B}_{\text{Đã biết}}
        \]
    \end{minipage}\hfill
    \begin{minipage}{.36\linewidth}
        \centering
        \begin{tikzpicture}[scale=.75]
            \draw[deepblue,thick] (-1.7,-1.3) rectangle (2.7,1.3);
            \fill[goldyellow!35] (1,0) circle(1);
            \begin{scope}
                \clip (0,0) circle(1);
                \fill[amberorange!70] (1,0) circle(1);

            \end{scope}
            \draw[deepblue] (0,0) circle(1);
            \draw[deepblue,thick] (1,0) circle(1);
            \node at (-.5,0){$A$};
            \node at (1.5,0){$B$};
            \node[below] at (.5,-1.35){Chỉ xét trong $B$};
        \end{tikzpicture}
    \end{minipage}
\end{theorybox}

% --- Câu 7 ---
\questionLinesManual{7}{7}{%
    Gieo một xúc xắc cân đối.
    Biết xuất hiện số chấm lẻ, tính xác suất xuất hiện mặt 1 chấm.
}{}

% --- Câu 8 ---
\questionLinesManual{7}{8}{%
    Một công ty có $P(A)=0{,}5$, $P(B)=0{,}6$, $P(A\cap B)=0{,}4$, với $A,B$ là thắng dự án 1, 2.
    Biết công ty không thắng dự án 1, tính xác suất thắng dự án 2.
}{}

\newpage

% =========================================================================
% TRANG 5: DẠNG 4
% =========================================================================
\dangbox{4}{Đếm số cách: chọn nhóm hay xếp thứ tự?}

\begin{theorybox}
    \[
        \begin{array}{c|c|c}
            \text{Chọn nhóm}&\text{Chọn và xếp thứ tự}&\text{Xếp tất cả}\\[7pt]
            \displaystyle C_n^k=\frac{n!}{k!(n-k)!}&\displaystyle A_n^k=\frac{n!}{(n-k)!}&P_n=n!
        \end{array}
    \]
    \[
        \boxed{P(A)=\frac{\text{Số cách chọn thỏa yêu cầu}}{\text{Tổng số cách chọn}}}\quad\text{(đồng khả năng).}
    \]
    \centering Đổi chỗ có tạo kết quả khác?\quad Không $\to C_n^k$;\quad Có $\to A_n^k$.
\end{theorybox}

% --- Câu 9 ---
\questionLinesManual{7}{9}{%
    Hộp có 7 bi trắng và 3 bi đỏ.
    Lấy ngẫu nhiên đồng thời 4 bi.
    Tính xác suất lấy được đúng 2 bi đỏ.
}{}

% --- Câu 10 ---
\questionLinesManual{7}{10}{%
    Vẫn hộp trên, lấy ngẫu nhiên đồng thời 4 bi.
    Tính xác suất lấy được ít nhất một bi đỏ.
}{}

\newpage

% =========================================================================
% TRANG 6: DẠNG 5
% =========================================================================
\dangbox{5}{Rút lần lượt: sau lần đầu còn lại gì?}

\begin{theorybox}
    \[
        \boxed{P(A\cap B)=P(A)\,P(B\mid A)}
    \]
    \begin{center}
        \begin{tikzpicture}[>=stealth]
            \node[draw=deepblue,rounded corners,inner sep=7pt] (a) {Ban đầu: $n$ vật};
            \node[draw=amberorange,rounded corners,inner sep=7pt,] (b) at (8,0) {Còn $n-1$ vật};
            \draw[->,thick,deepblue] (a)--node[above=5pt]{Không hoàn lại}(b);
        \end{tikzpicture}
    \end{center}
    \centering Có hoàn lại và xáo trộn $\Rightarrow$ số lượng trở về như cũ.
\end{theorybox}

% --- Câu 11 ---
\questionLinesManual{7}{11}{%
    Bình có 30 bi đen và 20 bi trắng.
    Lấy lần lượt 2 bi không hoàn lại.
    Tính xác suất lần đầu lấy bi đen, lần thứ hai lấy bi trắng.
}{}

% --- Câu 12 ---
\questionLinesManual{7}{12}{%
    Với bình trên, tính xác suất lấy được hai bi khác màu khi rút lần lượt 2 bi không hoàn lại.
}{}

\newpage

% =========================================================================
% TRANG 7: DẠNG 6
% =========================================================================
\dangbox{6}{Biết có ít nhất một: xác định đúng điều kiện}

\begin{theorybox}
    \[
        \text{Điều đã biết: }A\cup B\qquad\text{Điều cần tìm: }B.
    \]
    \[
        B\subset A\cup B\quad\Rightarrow\quad
        \boxed{P(B\mid A\cup B)=\frac{P(B)}{P(A\cup B)}}.
    \]
    \centering ``Có bi đỏ'' gồm cả trường hợp hai bi cùng đỏ.
\end{theorybox}

% --- Câu 13 ---
\questionLinesManual{7}{13}{%
    Hộp I có 6 bi đỏ, 4 bi xanh; hộp II có 5 bi đỏ, 5 bi xanh.
    Lấy độc lập một bi từ mỗi hộp.
    Biết có bi đỏ được lấy, tính xác suất bi lấy từ hộp II là đỏ.
}{}

% --- Câu 14 ---
\questionLinesManual{7}{14}{%
    Gieo hai xúc xắc cân đối một cách độc lập.
    Biết có ít nhất một con xuất hiện mặt 6, tính xác suất cả hai con đều xuất hiện mặt 6.
}{}

\newpage

% =========================================================================
% TRANG 8: DẠNG 7
% =========================================================================
\dangbox{7}{Hai điều kiện: chia trường hợp không trùng nhau}

\begin{theorybox}
    \[
        \text{Chọn nhóm}\ \longrightarrow\ n(\Omega)=C_n^k.
    \]
    \[
        \text{Các trường hợp rời nhau}\ \Rightarrow\ \boxed{n(A)=N_1+N_2+\cdots}.
    \]
    \centering Trong mỗi trường hợp: thực hiện nhiều bước $\Rightarrow$ nhân số cách.
\end{theorybox}

% --- Câu 15 ---
\questionLinesManual{7}{15}{%
    Có 3 giáo viên nữ Toán, 5 giáo viên nam Toán và 4 giáo viên nam Vật lý.
    Chọn ngẫu nhiên 3 người.
    Tính xác suất đoàn có đủ hai môn và đủ hai giới.
}{}

% --- Câu 16 ---
\questionLinesManual{7}{16}{%
    Hộp có 4 bi đỏ, 3 bi xanh và 2 bi vàng.
    Chọn ngẫu nhiên đồng thời 3 bi.
    Tính xác suất nhóm được chọn có đúng hai màu.
}{}

\newpage

% =========================================================================
% TRANG 9: DẠNG 8
% =========================================================================
\dangbox{8}{Hai xác suất có điều kiện: dùng bảng bốn ô}

\begin{theorybox}
    \[
        t=P(\overline A\cap\overline B),\quad x=P(A\cap\overline B),\quad y=P(\overline A\cap B).
    \]
    \[
        \boxed{P(A\mid\overline B)=\frac{x}{x+t}}\qquad
        \boxed{P(B\mid\overline A)=\frac{y}{y+t}}.
    \]
    \centering Tổng bốn ô $=1$;\qquad $P(A\cup B)=1-t$.
\end{theorybox}

% --- Câu 17 ---
\questionLinesManual{0}{17}{%
    Cho $P(A\mid\overline B)=3/5$, $P(B\mid\overline A)=2/3$ và $P(A\cap B)=1/10$.
    Tính $P(A\cup B)$.
}{}

\bangtrong
\drawdashlines{4}

% --- Câu 18 ---
\questionLinesManual{6}{18}{%
    Cho $P(A\mid\overline B)=2/5$, $P(B\mid\overline A)=4/7$ và $P(A\cap B)=0{,}1$.
    Tính $P(A\cup B)$.
}{}

\end{document}
